\documentclass[11pt,UTF8,fontset=fandol]{ctexart}
\usepackage{amsmath,amssymb,amsfonts}
\usepackage{graphicx,float,caption,subcaption}
\usepackage{booktabs,array,multirow,makecell}
\usepackage{geometry,enumitem,xcolor}
\usepackage[hidelinks]{hyperref}
\usepackage{url}
\geometry{margin=2.3cm}
\graphicspath{{/Users/yye/project/论文/assets/kgagent_guide/}}
\captionsetup{font=small,labelfont=bf}
\setlength{\parskip}{0.4em}
\hypersetup{pdftitle={KG-Agent 中文精读导读},pdfauthor={Jiang et al.}}

% 小白知识卡片：上下加线 + 粗体标签，稳健不挑内容
\newenvironment{tip}[1]{%
  \par\smallskip\noindent\rule{\linewidth}{0.4pt}\par\nobreak
  \noindent\textbf{【小白补丁】#1}\par\nobreak\noindent\ignorespaces
}{%
  \par\noindent\rule{\linewidth}{0.4pt}\par\smallskip
}

\title{KG-Agent 中文精读导读：让一个“小”模型自己在知识图谱上走多跳推理}
\author{原文：Jiang, Zhou, Zhao, Song, Zhu, Zhu, Wen（ACL 2025 / arXiv 2402.11163）}
\date{}

\begin{document}
\maketitle
\begin{center}\small\textit{本导读为面向初学者的中文学习笔记：含要点翻译、方法精讲、框架图解读与关键结果重述。}\end{center}

\section*{论文卡片}
\begin{center}
\begin{tabular}{@{}ll@{}}
\toprule
\textbf{标题} & KG-Agent: An Efficient Autonomous Agent Framework for \\
            & Complex Reasoning over Knowledge Graph \\
\textbf{作者} & Jinhao Jiang, Kun Zhou, Wayne Xin Zhao, Yang Song, \\
            & Chen Zhu, Hengshu Zhu, Ji-Rong Wen \\
\textbf{机构} & 中国人民大学高瓴人工智能学院、信息学院；BOSS 直聘 \\
\textbf{发表} & ACL 2025（arXiv 预印本 2402.11163，2024 年 2 月） \\
\textbf{方向} & 知识图谱（KG）$\times$ 大模型 $\times$ Agent $\times$ 指令微调 \\
\textbf{一句话} & 用约 1 万条“代码式”指令微调一个 7B 小模型，\\
              & 让它自己决定每一步调哪个工具、在知识图谱上多跳推理，\\
              & 超过用更大模型或更多数据的 SOTA 方法。 \\
\bottomrule
\end{tabular}
\end{center}

\section{一、这篇论文在解决什么问题}

大语言模型（LLM，如 GPT、LLaMA）光靠“脑子里记住的知识”去回答\textbf{复杂的事实性问题}时会力不从心，
尤其是\textbf{多跳（multi-hop）、知识密集型}的问题。比如“C 罗在 2011 年效力的球队里，哪一支成立得最晚？”——
这需要先查“C 罗 2011 年效力哪些队”，再查“每支队成立于哪一年”，最后“比大小取最晚的”，环环相扣。

\textbf{知识图谱（Knowledge Graph, KG）}恰好以“图”的方式存了海量事实，正好可以给 LLM 补充外部知识。
但 KG 数据量大、又是结构化的，LLM 不容易直接用好它。已有做法主要分两类：

\begin{enumerate}[leftmargin=2em]
  \item \textbf{检索增强（retrieval-augmented）}：先从 KG 里捞出相关三元组、拼成一段文字塞进提示里。
        缺点是丢掉了图的结构信息，还容易捞进一堆无关知识。
  \item \textbf{协同增强（synergy-augmented）}：设计一套“LLM 与 KG 反复交互”的机制，让 LLM 多次查询 KG。
        效果更好，但作者指出它有\textbf{两个硬伤}：
        \begin{itemize}[leftmargin=1.5em]
          \item \textbf{交互流程是“写死的”}：往往按人工设计好的固定多轮计划走，遇到难度或约束不同的新题型就不灵活；
          \item \textbf{依赖闭源大模型}：大多要靠 ChatGPT / GPT-4 这类强力闭源 API 去理解或学习解题，
                而它们蒸馏出来的“计划”未必适合指导更弱的小模型。
        \end{itemize}
\end{enumerate}

\textbf{本文的目标}：做一个\textbf{自主（autonomous）}的、基于 LLM 的 Agent 框架——叫 \textbf{KG-Agent}——
让一个\textbf{相对“小”的 7B 模型}能\textbf{自己做决定}，一步步在 KG 上推理直到找到答案，\textbf{全程不需要人类干预、也不依赖闭源 API}。

\section{二、摘要中译}

\begin{quote}
本文旨在提升大语言模型（LLM）在知识图谱（KG）上的推理能力，以回答复杂问题。受已有“设计 LLM 与 KG 交互策略”
工作的启发，我们提出一个基于 LLM 的自主 Agent 框架，称为 KG-Agent，它使一个小型 LLM 能够\textbf{主动做决策}，
直至完成在 KG 上的整个推理过程。在 KG-Agent 中，我们集成了 LLM、多功能工具箱（multifunctional toolbox）、
基于 KG 的执行器（KG-based executor）以及知识记忆（knowledge memory），并设计了一套\textbf{迭代机制}，
它能自主地选择工具、随后更新记忆，从而在 KG 上推理。为保证效果，我们利用程序语言（program language）来表述
KG 上的多跳推理过程，并合成一个\textbf{基于代码的指令数据集}来微调基座 LLM。大量实验表明：仅用 \textbf{1 万（10K）条样本}
微调 \textbf{LLaMA-7B}，就能在域内（in-domain）和域外（out-of-domain）数据集上\textbf{超过那些使用更大 LLM 或更多数据}
的当前最优方法。我们的代码与数据将公开发布。
\end{quote}

\section{三、三大核心贡献}
\begin{enumerate}[leftmargin=2em]
  \item \textbf{为 KG 量身定做一个工具箱}：把“操作结构化数据”的能力外挂给 LLM，让它能在 KG 与中间结果上
        做\textbf{筛选、计数、检索}等离散或高级操作（共三类工具，详见第七节）。
  \item \textbf{用现成数据“反向”合成代码式指令数据}：不去蒸馏 GPT-4，而是利用已有 KGQA 数据集里标注好的 SQL 查询，
        先把它“翻译”成 KG 推理程序（一串函数调用），再拆成多步的输入-输出对，做成指令微调数据。
  \item \textbf{提出自主迭代机制}：基于\textbf{“选工具 + 更新记忆”}的循环，把微调后的 LLM、工具箱、KG 执行器、
        知识记忆四者串起来，实现对 KG 的自主推理。
\end{enumerate}

\section{四、预备知识（小白先看这里）}

\begin{tip}{什么是知识图谱（Knowledge Graph, KG）}
KG 用“图”来存事实，基本单位是\textbf{三元组（triple）}$\langle e, r, e'\rangle$，意思是“头实体 $e$ 与尾实体 $e'$
之间存在关系 $r$”。例如 $\langle\text{C 罗}, \text{效力球队}, \text{葡萄牙国家队}\rangle$。很多三元组连起来就织成一张大图。
形式化地：$G=\{\langle e, r, e'\rangle \mid e, e'\in E,\ r\in R\}$，其中 $E$ 是实体集、$R$ 是关系集。
每个实体还有一个类型 $t$（如“国家”“人物”）。常见的大型 KG 有 Freebase、Wikidata。
\end{tip}

\begin{tip}{什么是多跳推理（multi-hop reasoning）}
答案离问题中提到的实体“隔了好几步”。比如 WebQSP 的答案一般在主题实体的 2 跳以内，CWQ 可达 4 跳。
每“跳”一步就相当于在图上沿一条关系走一格。多跳问题往往还带\textbf{约束}（如“=2011 年”）或\textbf{数值操作}（如“取最晚”）。
\end{tip}

\begin{tip}{什么是 Agent（智能体）与工具调用（tool calling）}
Agent 指让 LLM 不只“一次性吐答案”，而是像人一样\textbf{反复“观察—决策—行动”}：看一眼当前状态，
决定下一步干什么，调用一个外部“工具”（一段能执行的函数/程序），拿到结果后再继续。
代表性思路是 ReAct（把 LLM 变成会“边想边做”的语言智能体）。本文的工具就是一组“在 KG 上查询/筛选/计数”的函数。
\end{tip}

\begin{tip}{什么是指令微调（instruction tuning / SFT）}
拿“输入 $\to$ 期望输出”这样的成对样本去继续训练一个预训练好的 LLM，让它学会按指令办事。
本文的特别之处在于：输入是“问题 + 工具箱说明 + 当前已知的 KG 信息 + 历史程序”，输出是“下一步该调用的那个函数”——
于是模型学会的不是“背答案”，而是\textbf{“在每一步该调哪个工具”}这种决策能力。训练时只在“输出 $y$”上算交叉熵损失：
\end{tip}

\begin{equation*}
  \mathcal{L}=-\sum_{k=1}^{m}\log \Pr\big(y_k \mid x,\ y_{<k}\big). \tag{1}
\end{equation*}
其中 $x,y$ 是某一步的输入与标准答案，$m$ 是输出 $y$ 的 token 数，$y_k$、$y_{<k}$ 是第 $k$ 个及之前的 token。
直白说：给定输入和已经写出的部分，让模型预测下一个 token 尽量贴近标准答案。

\section{五、数据与任务：在哪些“考卷”上比}

KG-Agent 在两类任务上评测：\textbf{基于 KG 的问答（KGQA，域内）}和\textbf{开放域问答（ODQA，域外）}。

\begin{itemize}[leftmargin=2em]
  \item \textbf{域内（in-domain）}：用四个常用 KGQA 数据集做指令微调与测试——\textbf{WebQSP}（答案在 2 跳内）、
        \textbf{CWQ}（在 WebQSP 上加约束、更难，可达 4 跳）、\textbf{GrailQA}（专门考“泛化”，分 I.I.D./组合/零样本三档），
        前三者基于 \textbf{Freebase}；以及 \textbf{KQA Pro}（基于 \textbf{Wikidata}，需要组合推理、多跳、数值比较、集合运算等多种能力）。
  \item \textbf{域外（out-of-domain）}：微调后\textbf{不再训练}，直接零样本测三个 ODQA 数据集——\textbf{WQ}、\textbf{NQ}、\textbf{TQ}
        （分别记作 WQ-Freebase、NQ-Wiki、TQ-Wiki，即只保留能链接到 KG 的问题）。
  \item \textbf{换领域 KG}：还在 \textbf{MetaQA}（电影领域 KG，含 1/2/3 跳三个子集）上测“迁移到新 KG”的能力。
\end{itemize}

\begin{tip}{两个常用评测指标}
\textbf{Hits@1}：排在第一位的答案是否命中正确答案（看“头名准不准”）；
\textbf{F1}：综合考虑“查得全不全（召回）”和“查得准不准（精确）”，适合答案是一个集合的情况。两者都是越高越好。
\end{tip}

\textbf{关键设置}：指令微调数据\textbf{总共只采样 10{,}000 条}，按 \textbf{WebQSP : KQA Pro : GrailQA : CWQ $=$ 1 : 5 : 5 : 10}
的比例从域内数据集里抽取；基座模型用 \textbf{LLaMA2-7B}；学习率 2e-5、batch size 256、最大长度 1500、训练 3 个 epoch。
（推理过程默认假设问题里提到的实体已经给定。）

\section{六、方法逐节精讲}

KG-Agent 的核心是\textbf{一个被“调教好”的 LLM}，它能在 KG 上推理时\textbf{自主做决策}。整套方法分三块：
（4.1）为 KG 设计\textbf{工具箱}；（4.2）合成\textbf{代码式指令数据}并微调小模型；（4.3）设计\textbf{自主推理}框架。
总体框架见图~\ref{fig:arch}。

\begin{figure}[H]\centering
  \includegraphics[width=0.99\linewidth]{fig1_overview.pdf}
  \caption{KG-Agent 总体框架（原文 Figure 1，图内英文为原文保留）。\textbf{上半部分}是 Agent 的工作流：
  以问题“C 罗 2011 年效力的球队里哪支成立最晚？”为例，\emph{Planner（规划器，即微调后的 LLM）}根据\emph{知识记忆}
  在每一步选一个工具，\emph{Executor（执行器）}用程序编译器去 KG 上真正执行该工具调用，并把新得到的实体/关系
  写回\emph{知识记忆}（中间“Knowledge Memory Updating”一栏展示了 $t=1,2,3,\dots$ 各步记忆的变化），
  如此循环直到调用 \texttt{end} 结束（右侧“Autonomous Reasoning”）。\textbf{下半部分}展示了指令数据的合成：
  把标注 SQL 接地（grounding）到 KG 得到\emph{查询图（Query Graph）}，再转成一串\emph{KG 推理程序}
  （\texttt{get\_relation}、\texttt{get\_tail\_entity}、\texttt{get\_entity\_by\_constraint}$\dots$\texttt{end}），
  并按统一模板拆成多个“输入 $x_t$ / 输出 $y_t$”对。}
  \label{fig:arch}
\end{figure}

\subsection{4.1 为知识图谱设计工具箱}

LLM 不擅长精确操作结构化数据，于是作者为它配一套“工具”。按“在 KG 上推理通常需要的三种基本操作”——
\textbf{从 KG 取信息、按语义筛掉无关信息、对取出的信息做运算}——设计了\textbf{三类共若干工具}：

\begin{itemize}[leftmargin=2em]
  \item \textbf{抽取类（Extraction）}：访问 KG 本身。如取关系 \texttt{get\_relation}、取头/尾实体
        \texttt{get\_head\_entity}/\texttt{get\_tail\_entity}、按类型取实体 \texttt{get\_entity\_by\_type}、
        按约束取实体 \texttt{get\_entity\_by\_constraint}（支持 $=,>,\ge,<,\le$ 比较，或 \texttt{argmax}/\texttt{argmin} 取极值）。
  \item \textbf{逻辑类（Logic）}：对取出的实体集做运算。如计数 \texttt{count}、求交 \texttt{intersect}、求并 \texttt{union}、
        条件判断 \texttt{judge}，以及\textbf{结束推理并返回最终答案的} \texttt{end}。
  \item \textbf{语义类（Semantic）}：借助预训练模型实现的高级功能。如关系检索 \texttt{retrieve\_relation}
        （从候选关系里挑出与问题语义最相关的）和实体消歧 \texttt{disambiguate\_entity}。
\end{itemize}

\begin{tip}{为什么这样设计很灵活}
每个工具的“输入 $\to$ 输出”格式都按统一方式定义好了（例如多数工具的输入/输出都是“实体集合 $\{e\}$”），
所以工具箱可以\textbf{按真实需求随时扩展}，框架不挑具体 KG，也不挑任务类型。
\end{tip}

\subsection{4.2 合成代码式指令数据并微调小模型}

\textbf{思路}：不蒸馏 GPT-4，而是\textbf{“废物利用”现成 KGQA 数据集里标注好的 SQL 查询}。SQL 里本来就含有关系链、
条件、约束等信息，正好拿来合成“推理程序”。流程分三步：

\textbf{(1) 推理链抽取}：KG 太大，先用规则匹配从 KG 上拿到一小块与问题相关的\textbf{查询图（query graph）}——
它是树状的、能直接对应到逻辑形式，清晰刻画了“为得到答案，SQL 是怎么一步步执行的”。
再从问题里提到的实体（如“C 罗”）出发，用\textbf{广度优先搜索（BFS）}遍历查询图所有节点，
得到一条把起点实体连到答案实体的\textbf{推理链}（如 teams $\to$ roster\_team），途中自然带上约束（如 roster\_from = “2011”）
和数值操作（如 founded 取“最晚”）。

\textbf{(2) 推理程序生成}：把推理链转成若干\textbf{相互关联的三元组}（每个三元组≈一个中间推理步），
再把三元组改写成\textbf{代码格式的函数调用}（代表“调用某个工具”，且能真去 KG 上执行拿到对应三元组）。
具体地，从 \texttt{get\_relation(e)} 拿到当前候选关系，选一个关系传给 \texttt{get\_tail\_entity} 等函数得到新实体，
按推理链顺序生成所有函数调用，拼成完整的 \textbf{KG 推理程序}（见图~\ref{fig:arch} 下半部分的代码示例）。

\textbf{(3) 指令合成（输入-输出对构造）}：把这段程序里\textbf{每一个函数调用}都做成一条指令——
\textbf{输入} $=$ 问题 $+$ 工具箱定义 $+$ 当前 KG 信息（即当前实体集的下一批候选关系）$+$ 此步之前的历史程序；
\textbf{输出} $=$ 当前这一步的函数调用。执行完当前步后，“历史程序”和“当前 KG 信息”会相应更新，输出顺延为下一步的函数调用。
这样每条原始样本就被拆成\textbf{一串“输入-输出对”}，完整刻画了在 KG 上的推理轨迹。
对每条样本，把整条轨迹写成 $\{\langle x_1,y_1\rangle,\dots,\langle x_n,y_n\rangle\}$（$n$ 为总步数），
喂给\textbf{仅解码器（decoder-only）}的 LLaMA-7B 做 SFT，损失见公式 (1)。

\subsection{4.3 自主推理：四个组件 + “选工具/更记忆”循环}

微调完之后，作者把它装进一个 Agent 框架（图~\ref{fig:arch} 上半部分），含\textbf{四个组件}：

\begin{enumerate}[leftmargin=2em]
  \item \textbf{LLM 规划器（Planner）}：就是 4.2 微调出的那个 LLM，负责“看记忆、选工具”。
  \item \textbf{多功能工具箱（Toolbox）}：4.1 里那套工具。
  \item \textbf{KG 执行器（Executor）}：用程序编译器真正执行函数调用，能缓存/操作中间变量、从 KG 抽出新实体或关系。
  \item \textbf{知识记忆（Knowledge Memory）}：记录全程的上下文与“当前有用信息”。
\end{enumerate}

\textbf{知识记忆的四部分}：①自然语言问题、②工具箱定义、③当前 KG 信息、④历史推理程序。
前两部分用给定的问题和工具箱初始化，\textbf{全程不变}；后两部分初始化为空列表，\textbf{每一步都更新}。

整个推理是一个\textbf{“选工具 $\to$ 执行 $\to$ 更新记忆”的循环}：

\begin{itemize}[leftmargin=2em]
  \item \textbf{规划器选工具}：把当前记忆按模板拼成输入，LLM 从中\textbf{生成一个函数调用}（选一个工具 + 它的参数）。
        它要应对四类需求：把提到的实体链接到 KG、访问 KG 信息、处理中间结果、或调用 \texttt{end} 返回最终答案。
  \item \textbf{执行器更记忆}：执行该函数调用后更新记忆——把这次调用追加进“历史程序”；若该工具是去 KG 取新信息
        （如 \texttt{get\_relation}），就把结果加进“当前 KG 信息”。
  \item \textbf{自主迭代}：不断重复上面两步做逐步推理。作者打了个比方：这个多轮决策过程\textbf{就像“沿着关系在 KG 上行走”}。
        一旦走到答案实体，Agent 会\textbf{自动停止}。整个过程\textbf{不依赖具体任务类型或具体 KG}，因此是个通用框架。
\end{itemize}

\begin{tip}{和已有工作的关键区别（对照原文 Table 1）}
原文用一张表对比了 Pangu、StructGPT、RoG、ChatDB、KB-BINDER 等方法。KG-Agent 的与众不同在于同时满足：
交互方式是\textbf{自主（auto）}而非写死（pd）、基座是\textbf{开源 LLaMA2-7B}（不靠闭源 API）、
\textbf{同时具备工具（Tool）与记忆（Memory）}、并支持\textbf{跨多个 KG 的多任务（Multi-Task）泛化}——这一整套组合是它独有的。
\end{tip}

\section{七、实验与结果：1 万样本的 7B 如何超过 SOTA}

\textbf{对比对象}：域内比三类基线——子图推理（如 NSM、UniKGQA）、LM 序列生成（如 RNG-KBQA、TIARA、PanGu）、
LLM 方法（如 RoG、StructGPT，以及 ChatGPT/Davinci-003/GPT-4）；域外比微调类（T5/BART）和 LLM 类。

\textbf{域内主结果（基于 Freebase 的 WebQSP/CWQ/GrailQA，数值源自原文 Table 2，节选关键行；越高越好，加粗为最优）}：

\begin{table}[H]\centering
\caption{基于 Freebase 的域内主结果对比（数值源自原文 Table 2）。GrailQA 报告 F1。}
\begin{tabular}{lcccccc}
\toprule
\multirow{2}{*}{\textbf{方法}} & \multicolumn{2}{c}{\textbf{WebQSP}} & \multicolumn{2}{c}{\textbf{CWQ}} & \multicolumn{2}{c}{\textbf{GrailQA (F1)}} \\
\cmidrule(lr){2-3}\cmidrule(lr){4-5}\cmidrule(lr){6-7}
 & Hits@1 & F1 & Hits@1 & F1 & Overall & Zero-shot \\
\midrule
UniKGQA（子图推理）   & 75.1 & 70.2 & 50.7 & 48.0 & --   & --   \\
RoG（LLM 方法）       & 85.7 & 70.8 & 62.6 & 56.2 & --   & --   \\
FC-KBQA（LM 生成）    & --   & 76.9 & --   & 56.4 & 83.8 & 83.1 \\
PanGu w/ T5-3B        & --   & 79.6 & --   & --   & 83.4 & --   \\
GPT-4                 & 73.2 & 62.3 & 55.6 & 49.9 & 31.7 & 39.2 \\
StructGPT             & 72.6 & 63.7 & 54.3 & 49.6 & 54.6 & 50.5 \\
\textbf{KG-Agent（本文）} & \textbf{83.3} & \textbf{81.0} & \textbf{72.2} & \textbf{69.8} & \textbf{86.1} & \textbf{86.3} \\
\bottomrule
\end{tabular}
\end{table}

\textbf{怎么读这张表}：KG-Agent 在 F1 上几乎全面领先——WebQSP 的 F1 达 81.0（高于 PanGu-3B 的 79.6），
CWQ 的 F1 达 69.8、GrailQA 的 Overall F1 达 86.1，都明显超过此前最优。直接用 GPT-4/ChatGPT 与最好的微调方法
仍有较大差距，说明“光靠 LLM 硬答复杂问题”很难。

\begin{tip}{关键数字：少数据、小模型，却更强}
\textbf{只用约 10K 条样本}微调一个\textbf{7B}小模型，就超过了那些用\textbf{更大 LLM 或全量数据}的方法。
原文给出更直观的对比：相比基线，KG-Agent 在 CWQ 和 GrailQA 上分别只用了约 \textbf{36\%} 和 \textbf{23\%} 的原始训练集量，
却分别取得 \textbf{7.5\%} 和 \textbf{2.7\%} 的 F1 相对提升；在混合数据上联合训练（多数据集相互增益）后，
WebQSP / CWQ / GrailQA 的 F1 又分别再涨 \textbf{1.7\% / 7.5\% / 2.7\%}。
\end{tip}

\textbf{基于 Wikidata 的 KQA Pro（数值源自原文 Table 3，准确率，越高越好）}：

\begin{table}[H]\centering
\caption{KQA Pro 测试集准确率（源自原文 Table 3，节选）。Multi-hop/Comparison/Count 等为按推理类型细分的子集。}
\begin{tabular}{lcccccc}
\toprule
\textbf{方法} & Overall & Multi-hop & Comparison & Logical & Count & Zero-shot \\
\midrule
BART SPARQL  & 89.68 & 88.49 & 96.12 & 88.67 & 85.78 & 87.88 \\
GPT-4        & 37.43 & 34.82 & 55.75 & 36.81 & 15.27 & 27.28 \\
ChatGPT      & 24.96 & 24.22 & 39.15 & 25.51 & 10.76 & 15.67 \\
\textbf{KG-Agent（本文）} & \textbf{92.15} & \textbf{91.03} & \textbf{96.32} & \textbf{91.28} & \textbf{88.21} & \textbf{91.40} \\
\bottomrule
\end{tabular}
\end{table}

KG-Agent 在 KQA Pro 总体准确率达 \textbf{92.15}，全面超过强力的 BART SPARQL（89.68）；而 GPT-4/ChatGPT 在这类
需要组合与数值推理的题上表现很差（如 Count 子集仅 15.27 / 10.76），凸显“配上工具与 KG”的价值。

\textbf{域外零样本（源自原文 Table 4）}：微调后\textbf{不再训练}，直接零样本测 ODQA。在基于 Freebase 的 WQ-Freebase 上，
T5/BART 即使全量微调也答不好；KG-Agent 因为“只需学会与 KG 交互、不必死记具体知识”，
零样本就比这些全量监督模型更好（原文摘要：WQ-Freebase 与 TQ-Wiki 上准确率分别相对提升约 \textbf{9.7\%} 和 \textbf{8.5\%}）。

\textbf{迁移到领域 KG（源自原文 Table 5，MetaQA 电影 KG）}：在 1/2/3 跳子集上，KG-Agent（97.1 / 98.0 / 92.1）
一致超过 StructGPT 与监督微调基线，说明它学到的是\textbf{“在 KG 上推理的通用能力”}，能高效迁移到新 KG。

\textbf{进一步分析}：
\begin{itemize}[leftmargin=2em]
  \item \textbf{指令数量的影响（原文 Figure 2）}：把总量从 2k 指数级加到 64k，性能随数据增多而上升、最终趋于稳定；
        从 16k 到 64k 提升已不明显——作者认为这与数据“多样性”有关，未来会构造更多样的样本。这也佐证了
        \textbf{“约 10K 已足够”}的结论。
  \item \textbf{采样比例的影响（原文 Table 6）}：在总量不变下，把某个数据集的比例翻倍，它自己的分数会涨，
        但\textbf{三者平均分}都不如作者选定的 1:5:5:10 比例——说明该比例较好地平衡了综合能力。
\end{itemize}

\section{八、小白知识点卡片（速查）}
\begin{itemize}[leftmargin=2em]
  \item \textbf{知识图谱 / KG}：用三元组 $\langle e,r,e'\rangle$ 存事实的图（Freebase、Wikidata 等）。
  \item \textbf{多跳推理}：答案离问题实体隔多步，常带约束与数值操作。
  \item \textbf{Agent / 工具调用}：让 LLM “观察—决策—行动”循环，调用可执行的外部函数。
  \item \textbf{指令微调（SFT）}：用“输入$\to$输出”成对样本继续训练 LLM；本文教模型学“每步该调哪个工具”。
  \item \textbf{KGQA / ODQA}：基于 KG 的问答 / 开放域问答（本文域内 / 域外任务）。
  \item \textbf{查询图（query graph）}：从 SQL 接地到 KG 得到的树状小图，刻画求解的执行流。
  \item \textbf{知识记忆}：问题 + 工具箱定义 + 当前 KG 信息 + 历史程序，后两者每步更新。
  \item \textbf{Hits@1 / F1}：头名是否命中 / 综合查准查全；越高越好。
  \item \textbf{SPARQL / SQL}：在图数据库 / 关系库上查询的结构化语言（本文借标注 SQL 反向合成程序）。
\end{itemize}

\section{九、一句话总结与谱系定位}

\textbf{一句话}：KG-Agent 把“在 KG 上做多跳推理”改写成\textbf{一段可执行的代码式轨迹}——
先用现成 KGQA 数据集的标注 SQL\textbf{反向合成}“代码式指令数据”，用约 \textbf{1 万条}就把一个 \textbf{7B 小模型}调成“规划器”；
再让它在\textbf{“选工具 $\to$ 执行 $\to$ 更新知识记忆”}的循环里\textbf{自己决定每一步}，像“沿关系在图上行走”一样找到答案，
从而\textbf{又准、又省、又不依赖闭源大模型}地超过更大/更多数据的 SOTA。

\textbf{在“LLM $\times$ 知识图谱推理”谱系里的位置}：它属于\textbf{“协同增强 + 自主 Agent + 小模型指令微调”}这一路线的代表作。
与“写死交互流程”的 StructGPT/KB-BINDER/RoG 相比，它的\textbf{交互流程是模型自主决定}的；
与依赖 GPT-4 的 ChatDB 等相比，它\textbf{只靠一个开源 7B}并\textbf{额外配了工具去做 count 等专门操作}；
与“先生成完整计划再接地”的方法相比，它\textbf{边走边决策}、灵活适配不同难度与约束。
你可以把它和本系列里走“图检索 + 软提示”的 G-Retriever、走“图上思维链”的 GraphCoT、走“强化学习”的方法\textbf{对照着看}——
它们各自把“图的信息”用进 LLM 的方式很不一样。

\textbf{动手提示}：复现的关键三步是——（1）从标注 SQL 合成 KG 推理程序与输入-输出对；
（2）用 1:5:5:10 比例约 10K 样本微调 LLaMA2-7B；（3）实现“规划器 + 工具箱 + KG 执行器 + 知识记忆”的迭代循环。
论文称代码与数据将公开发布。

\end{document}
